\begin{resumo}
Os rolamentos dos rodeiros ferroviários (\textit{axlebox bearings}) operam sob condições dinâmicas severas de carregamento e vibração, constituindo componentes críticos para a segurança e a regularidade do transporte ferroviário de cargas e passageiros. Falhas mecânicas não detectadas nesses elementos podem evoluir para o travamento de eixos e descolamentos estruturais, culminando em descarrilamentos com elevados custos operacionais e riscos à vida humana. Nesse contexto, a manutenção baseada na condição (\textit{Condition-Based Maintenance} -- CBM), fundamentada na análise de vibrações e no Aprendizado de Máquina (\textit{Machine Learning}), surge como uma alternativa preditiva superior às inspeções periódicas preventivas e aos detectores térmicos de via (\textit{hotbox detectors}), os quais identificam anomalias apenas em estágios terminais de degradação. Este trabalho apresenta o desenvolvimento de um \textit{framework} preditivo para diagnóstico e classificação multiclasse de defeitos em rolamentos mecânicos com base no banco de dados experimental de referência da \textit{Case Western Reserve University} (CWRU). O pipeline metodológico compreende a segmentação de sinais de aceleração adquiridos a \SI{12}{\kilo\hertz} em janelas discretas de \num{2048} pontos com preservação de linhagem temporal (\textit{data lineage}), a extração de atributos estocásticos no domínio do tempo (RMS, Curtose, Fator de Crista, Desvio Padrão, Assimetria, Pico a Pico, Fator de Forma e Fator de Impulso) e atributos espectrais baseados na Transformada Rápida de Fourier (FFT com janela de Hanning) e na Análise de Envelope via Transformada de Hilbert, monitorando as frequências características de defeito (BPFO, BPFI, BSF e FTF). Para a modelagem supervisionada, avaliaram-se os algoritmos \textit{Random Forest} e \textit{Support Vector Machines} (SVM com \textit{kernel} RBF e padronização \textit{StandardScaler}). Para evitar vazamento temporal de dados (\textit{data leakage}), empregou-se divisão estratificada por arquivo de ensaio (70/30) e validação cruzada estratificada agrupada (\textit{Stratified Group K-Fold}, $k=4$). Ambos os classificadores alcançaram acurácias superiores a \SI{99}{\percent} no conjunto de teste independente (RF com \SI{99,32}{\percent} e SVM com \SI{99,10}{\percent}). Em nível de ensaio físico completo, a agregação por votação majoritária (\textit{file-level majority voting}) atingiu \SI{100,00}{\percent} de acerto em todos os mancais avaliados, evidenciando a robustez da abordagem para aplicação em sistemas mecatrônicos de monitoramento ferroviário.

\vspace{\onelineskip}
\noindent
\textbf{Palavras-chave:} Manutenção Preditiva. Rolamentos Ferroviários. Processamento de Sinais. Machine Learning. CWRU Dataset.
\end{resumo}

\begin{resumo}[Abstract]
\begin{otherlanguage*}{english}
Railway axlebox bearings operate under severe dynamic loads, impacts, and temperature fluctuations, representing critical components for operational safety and reliability in rail transportation. Undetected mechanical defects in these bearings may progress to axle seizure, wheel derailment, severe financial losses, and hazardous accidents. In this context, Condition-Based Maintenance (CBM) leveraging vibration analysis and Machine Learning algorithms provides an advanced predictive alternative to conventional time-based maintenance and wayside infrared hotbox detectors, which only detect anomalies in terminal failure stages. This work develops an end-to-end computational framework for multi-class fault classification in rolling element bearings using the benchmark dataset from Case Western Reserve University (CWRU). The methodology encompasses signal segmentation into discrete \num{2048}-sample windows with temporal data lineage tracking, followed by feature extraction in both time domain (RMS, Kurtosis, Crest Factor, Standard Deviation, Skewness, Peak-to-Peak, Shape Factor, and Impulse Factor) and frequency domain (Fast Fourier Transform with Hanning windowing and Hilbert envelope spectrum analysis focused on characteristic fault frequencies: BPFO, BPFI, BSF, and FTF). Supervised classification models, specifically Random Forest and Support Vector Machines (SVM with RBF kernel and StandardScaler), were trained and validated. To strictly prevent temporal data leakage, an independent 70/30 train/test split at the file/experiment level and a 4-fold Stratified Group Cross-Validation were implemented. Both classifiers achieved generalization accuracies exceeding \SI{99}{\percent} on the test set (\SI{99.32}{\percent} for Random Forest and \SI{99.10}{\percent} for SVM). Furthermore, an experiment-level majority voting scheme achieved \SI{100.00}{\percent} diagnostic accuracy across all tested bearings, demonstrating the effectiveness and readiness of the proposed methodology for integration into intelligent mechatronic condition monitoring systems.

\vspace{\onelineskip}
\noindent
\textbf{Keywords:} Predictive Maintenance. Railway Bearings. Signal Processing. Machine Learning. CWRU Dataset.
\end{otherlanguage*}
\end{resumo}
